\section{Vetores e matrizes}

\begin{frame}[fragile,t]{11. Vetores}
O que é? Em fundamentos de programação, um vetor é uma estrutura unidimensional, organizada em posições sequenciais. Em Python introdutório, uma lista pode representar esse tipo de estrutura.\par
\begin{visibleenv}<2->
\EMASCcodelabel{Estrutura do comando}
\begin{lstlisting}
vetor = [elemento1, elemento2, elemento3]
\end{lstlisting}
\end{visibleenv}
\end{frame}

\begin{frame}[fragile,t]{11. Vetores / Exemplo}
\EMASCcodelabel{Exemplo prático}
\begin{lstlisting}
numeros = [10, 20, 30, 40, 50]

for i in range(len(numeros)):
    print(numeros[i])
\end{lstlisting}
\begin{visibleenv}<2->
\begin{EMASCnote}
Entenda o exemplo: \texttt{len(numeros)} informa a quantidade de elementos. \texttt{range()} gera os índices e \texttt{numeros[i]} acessa o elemento na posição indicada.
\end{EMASCnote}
\end{visibleenv}
\end{frame}

\begin{frame}[fragile,t]{11. Vetores / Acessando posições específicas}
\EMASCcodelabel{Exemplo prático}
\begin{lstlisting}
numeros = [10, 20, 30, 40, 50]

print(numeros[0])   # 10
print(numeros[2])   # 30
print(numeros[4])   # 50
\end{lstlisting}
\end{frame}

\begin{frame}[fragile,t]{12. Matrizes}
O que é? Matrizes organizam dados em linhas e colunas. Uma matriz bidimensional pode ser representada em Python por uma lista cujos elementos também são listas.\par
\begin{visibleenv}<2->
\EMASCcodelabel{Estrutura do comando}
\begin{lstlisting}
matriz = [
    [elemento1, elemento2],
    [elemento3, elemento4]
]
\end{lstlisting}
\end{visibleenv}
\end{frame}

\begin{frame}[fragile,t]{12. Matrizes / Exemplo}
\EMASCcodelabel{Exemplo prático}
\begin{lstlisting}
matriz = [
    [1, 2, 3],
    [4, 5, 6],
    [7, 8, 9]
]

print(matriz[0][0])    # 1
print(matriz[1][2])    # 6
print(matriz[2][1])    # 8
\end{lstlisting}
\begin{visibleenv}<2->
\begin{EMASCnote}
Entenda o exemplo: O primeiro índice indica a linha e o segundo indica a coluna. Como os índices começam em zero, \texttt{matriz[1][2]} acessa a segunda linha e a terceira coluna.
\end{EMASCnote}
\end{visibleenv}
\end{frame}

\begin{frame}[fragile,t]{12. Matrizes / Percorrendo todos os elementos}
\EMASCcodelabel{Exemplo prático}
\begin{lstlisting}
for linha in matriz:
    for elemento in linha:
        print(elemento)
\end{lstlisting}
\begin{visibleenv}<2->
\begin{EMASCnote}
O primeiro \texttt{for} percorre cada linha. O segundo percorre os elementos da linha atual.
\end{EMASCnote}
\end{visibleenv}
\end{frame}

\begin{frame}[t]{13. Resumo geral}
\footnotesize
\renewcommand{\arraystretch}{0.95}
\begin{tabularx}{\linewidth}{@{}lX@{}}
\toprule
\textbf{Assunto} & \textbf{Comando / representação} \\
\midrule
Variáveis & \texttt{nome = valor} \\
Tipos & \texttt{int, float, str, bool} \\
Entrada & \texttt{input()} \\
Saída & \texttt{print()} \\
Operadores & \texttt{+ - * / // \% **; == != \textgreater{} \textless{} \textgreater{}= \textless{}=} \\
Condicionais & \texttt{if, elif, else} \\
Repetição & \texttt{while, for} \\
Funções & \texttt{def, return} \\
Listas & \texttt{[]} \\
Tuplas & \texttt{()} \\
Dicionários & \texttt{\{\} com pares chave: valor} \\
Vetores & lista unidimensional \\
Matrizes & listas aninhadas (linhas e colunas) \\
\bottomrule
\end{tabularx}
\end{frame}
